The greedy allocation used by all expected-value policies is exactly
optimal for the per-day problem it solves.

\begin{proposition}
Let $D_1,\dots,D_S$ be non-negative random variables, $\omega_s > 0$,
$\kappa > 0$, and consider
\[
  \max_{x \in \mathbb{Z}_{\ge 0}^S, \; \sum_s x_s = B}
  \; \sum_{s=1}^S \omega_s\, \mathbb{E}\!\left[\min(D_s, \kappa x_s)\right].
\]
The greedy algorithm that assigns crews one at a time to the family with
the largest marginal gain attains the optimum.
\end{proposition}

\begin{proof}
For fixed $s$, define $g_s(x) = \omega_s \mathbb{E}[\min(D_s, \kappa x)]$
on integers $x \ge 0$. For any realization $d \ge 0$ the map
$c \mapsto \min(d, c)$ is concave in $c$; expectation preserves concavity,
and restriction to the integer grid yields non-increasing increments
$\Delta_s(x) = g_s(x+1) - g_s(x)$. The problem is therefore the
maximization of a separable discrete-concave function over the simplex
$\sum_s x_s = B$, for which the greedy/incremental algorithm is optimal:
at every step the chosen increment is at least as large as any increment
available later (by monotonicity of $\Delta_s$), so an exchange argument
converts any optimal solution into the greedy one without loss of
objective value.
\end{proof}

In the implementation $\mathbb{E}[\min(D_s, c)]$ is evaluated exactly for
the discrete distribution defined by linear interpolation of the inverse
CDF through the five forecast quantiles on a fixed 99-point grid
(\texttt{src/optimization/allocation.py}); point forecasts use a
degenerate distribution, and the oracle uses the realized value. The unit
test \texttt{test\_greedy\_matches\_bruteforce\_on\_small\_instance}
verifies greedy-vs-brute-force equality on random instances.
